% Informe de la práctica 4. La portada se agrega por separado.
\documentclass[12pt,a4paper]{article}
% ============================================================================
% Formato.tex: configuración común de los informes del Laboratorio de
% Circuitos Eléctricos II (LIE-311). Se incluye desde Documento.tex.
% La portada se omite intencionalmente y se agrega por separado.
% ============================================================================

% Codificación, idioma y tipografía. inputenc mantiene compatibilidad con
% pdfLaTeX; en LuaLaTeX/XeLaTeX puede omitirse.
\usepackage[utf8]{inputenc}
\usepackage[T1]{fontenc}
\usepackage[spanish,es-nodecimaldot]{babel}
\usepackage{lmodern}

% Matemáticas, unidades, tablas e imágenes.
\usepackage{amsmath,amssymb}
\usepackage{siunitx}                 % Ej.: \SI{10}{\volt}
\usepackage{graphicx}
\usepackage{booktabs}
\usepackage{array}
\usepackage{enumitem}
\usepackage{float}
\usepackage{caption}
\usepackage{xcolor}
\usepackage[american]{circuitikz}

% Diseño de página, separación de párrafos y enlaces.
\usepackage{geometry}
\usepackage{fancyhdr}
\usepackage{parskip}
\usepackage{csquotes}
\usepackage[hidelinks]{hyperref}

% El estilo numérico produce citas compactas, apropiadas para IEEE. Para APA,
% cambie style=numeric-comp por style=apa (y agregue el paquete correspondiente
% si su instalación de biblatex no lo incluye).
\usepackage[
    backend=biber,
    style=numeric-comp,
    sorting=none
]{biblatex}

% La comprobación permite compilar la plantilla antes de crear references.bib.
% Al entregar, ese archivo debe contener el manual y todas las fuentes usadas.
\IfFileExists{references.bib}{\addbibresource{references.bib}}{}

\geometry{a4paper, margin=2.5cm, headheight=15pt}
\setlength{\parindent}{0pt}
\setlength{\parskip}{6pt}
\setlist{itemsep=2pt, topsep=3pt}
\captionsetup{font=small, labelfont=bf}
\sisetup{locale=DE, per-mode=symbol, detect-all}

% Datos editables. Se redefinen en Documento.tex con \renewcommand. Se declaran
% antes del encabezado porque este último los utiliza.
\newcommand{\numeroPractica}{N}
\newcommand{\nombrepractica}{Título de la práctica}

% Encabezado y pie. Los valores de la práctica se definen en Documento.tex.
\pagestyle{fancy}
\fancyhf{}
\lhead{LIE-311 -- Laboratorio de Circuitos Eléctricos II}
\rhead{Práctica \numeroPractica}
\cfoot{\thepage}

% Inserta una figura centrada y numerada.
% Uso: \imagen[etiqueta]{archivo}{pie de figura}
% Las capturas virtuales deben mostrar fecha y hora; el archivo puede ser PNG,
% JPG o PDF. Una etiqueta permite referenciar la figura con \ref{fig:montaje}.
\newcommand{\imagen}[3][]{%
    \begin{figure}[H]
        \centering
        \includegraphics[width=0.88\textwidth]{#2}
        \caption{#3}
        \if\relax\detokenize{#1}\relax\else\label{#1}\fi
    \end{figure}%
}

% Entorno reutilizable para objetivos y conclusiones numerados.
\newenvironment{itemsinforme}
    {\begin{enumerate}[label=\arabic*., leftmargin=*]}
    {\end{enumerate}}

% Marcador visible para localizar contenido pendiente. Debe reemplazarse o
% eliminarse antes de entregar el informe final.
\newcommand{\rellenar}[1]{\textcolor{red}{\textbf{[Completar: #1]}}}

\hypersetup{
    pdftitle={Informe de laboratorio -- \nombrepractica},
    pdfauthor={Nombre del estudiante},
    pdfsubject={Laboratorio de Circuitos Eléctricos II -- LIE-311}
}

\usepackage{circuitikz}
\renewcommand{\numeroPractica}{4}
\renewcommand{\nombrepractica}{Diagramas de Bode y filtros}

\begin{document}

\section{Introducción}
La práctica caracterizó la respuesta en frecuencia de varios circuitos pasivos mediante mediciones y diagramas de Bode. La función de transferencia relacionó la señal de salida con la señal de entrada y permitió estudiar cómo la ganancia y la fase variaron cuando cambió la frecuencia. Esta representación resultó útil para reconocer las regiones de paso, atenuación y transición de un filtro. Se analizaron un circuito RC con salida en el capacitor, un circuito RL con salida en el inductor y dos configuraciones RLC con comportamiento pasa-banda y rechaza-banda. Para cada montaje se consideraron las frecuencias de corte, la ganancia expresada en decibelios y las diferencias entre los resultados teóricos, simulados y experimentales. El análisis permitió relacionar el comportamiento de los componentes reactivos con aplicaciones de selección y rechazo de frecuencias en sistemas eléctricos y de comunicación.

\section{Objetivos}
\begin{itemsinforme}
    \item Estudiar los diagramas de Bode y su importancia en el análisis de la respuesta en frecuencia de sistemas.
    \item Obtener y representar la magnitud y la fase de una función de transferencia.
    \item Analizar el comportamiento de filtros pasivos pasa-bajos, pasa-altos, pasa-banda y rechaza-banda.
    \item Comparar las ganancias obtenidas mediante medición, simulación y cálculo teórico.
\end{itemsinforme}

\section{Materiales y equipo}
\begin{itemize}
    \item Software Multisim.
    \item Generador de funciones.
    \item Osciloscopio.
    \item Multímetro digital.
    \item Bode Plotter.
    \item Resistencias.
    \item Capacitores.
    \item Inductores.
\end{itemize}

\section{Marco teórico}
\subsection{Función de transferencia}
La función de transferencia ``\(H(\omega)\)'', también llamada función de red, fue una herramienta analítica útil para determinar la respuesta a la frecuencia de un circuito. La respuesta a la frecuencia de un circuito fue la gráfica de la función de transferencia de este mismo \(H(\omega)\), donde \(\omega\) varió desde \(\omega=0\) hasta \(\omega=\infty\).

Una función de transferencia fue la relación entre una función forzada \(Y(\omega)\) y una función de excitación \(X(\omega)\), o entre una salida y una entrada, dependiente de la frecuencia. Su expresión fue
\begin{equation}
    H(\omega)=\frac{Y(\omega)}{X(\omega)}.
    \label{eq:transferencia}
\end{equation}

Suponiendo las condiciones iniciales iguales a cero, puesto que la entrada y salida pudieron ser una señal de tensión o de corriente en cualquier parte del circuito, existieron cuatro posibles funciones de transferencia: ganancia de voltaje \(H(\omega)=V_{out}(\omega)/V_{in}(\omega)\), ganancia de corriente \(H(\omega)=I_{out}(\omega)/I_{in}(\omega)\), transferencia de impedancia \(H(\omega)=V_{out}(\omega)/I_{in}(\omega)\) y transferencia de admitancia \(H(\omega)=I_{out}(\omega)/V_{in}(\omega)\).

Para obtener la función de transferencia, primero se obtuvo el equivalente en el dominio de la frecuencia del circuito, sustituyendo las resistencias, inductores y capacitores por los circuitos transformados correspondientes en función de la frecuencia. Después se utilizó cualquier técnica de circuitos para obtener la función de transferencia a analizar.

La función de transferencia pudo expresarse en términos de sus polinomios numerador \(N(\omega)\) y denominador \(D(\omega)\):
\begin{equation}
    H(\omega)=\frac{N(\omega)}{D(\omega)}.
    \label{eq:numeradordenominador}
\end{equation}
Las raíces de \(N(\omega)\) fueron los ceros de la función de transferencia y las raíces de \(D(\omega)\) fueron los polos. Un cero, como raíz del polinomio del numerador, fue un valor que produjo un valor cero de la función. Un polo, como raíz del polinomio del denominador, fue un valor para el cual la función fue infinita.

En el dominio de la frecuencia, la resistencia permaneció como \(R\), el inductor se representó mediante \(j\omega L\) y el capacitor mediante \(1/(j\omega C)\). Con estas impedancias se obtuvo la función de transferencia aplicando las leyes de circuitos.

En el dominio de la frecuencia, la resistencia permaneció como \(R\), el inductor se representó mediante \(j\omega L\) y el capacitor mediante \(1/(j\omega C)\). Con estas impedancias se obtuvo la función de transferencia aplicando las leyes de circuitos.

\subsection{Decibelios y diagramas de Bode}
Una forma más sistemática de obtener la respuesta a la frecuencia consistió en utilizar los diagramas de Bode. Estos se basaron en logaritmos y decibelios para expresar la ganancia. Las propiedades empleadas fueron \(\log(P_1P_2)=\log P_1+\log P_2\), \(\log(P_1/P_2)=\log P_1-\log P_2\), \(\log(P^n)=n\log P\) y \(\log 1=0\).

En los sistemas de comunicación, la ganancia se midió en bels. El bel se utilizó para medir la relación entre dos niveles de potencia o la ganancia de potencia \(G\), mientras que el decibelio correspondió a una décima parte de un bel:
\begin{equation}
    G_{\mathrm{dB}}=10\log_{10}\left(\frac{P_2}{P_1}\right).
    \label{eq:potenciadb}
\end{equation}
De manera alterna, la ganancia pudo expresarse en términos de la relación entre tensiones o corrientes. Para resistencias iguales se obtuvo la ganancia de tensión y de corriente como
\begin{equation}
    G_{\mathrm{dB}}=20\log_{10}\left|\frac{V_{\mathrm{o}}}{V_{\mathrm{i}}}\right|.
    \label{eq:db}
\end{equation}
También se pudo calcular la ganancia de corriente mediante \(G_{\mathrm{dB}}=20\log_{10}|I_2/I_1|\). Se utilizó el factor 20 para tensión o corriente debido a la relación al cuadrado entre potencia y dichas variables. El valor en dB fue una medición logarítmica de la relación entre dos variables del mismo tipo; por ello se aplicó a la ganancia de voltaje o corriente, pero no directamente a las expresiones de admitancia o impedancia de transferencia. Solo se utilizaron las magnitudes de la tensión y la corriente.

\subsection{Filtros y frecuencia de corte}
La respuesta a la frecuencia de un circuito fue la variación de su comportamiento al cambiar la frecuencia de la señal. El análisis de la respuesta a la frecuencia tuvo gran importancia en sistemas de comunicaciones y de control. Una aplicación fue el uso de filtros eléctricos, que bloquearon o eliminaron señales a frecuencias no deseadas y dejaron pasar señales a frecuencias deseadas.

Los diagramas de Bode fueron gráficas semilogarítmicas de la magnitud y de la fase de una función de transferencia en función de la frecuencia. La magnitud en decibelios se graficó contra el logaritmo de la frecuencia y, en un diagrama aparte, la fase en grados contra el logaritmo de la frecuencia. Al elaborar un diagrama de Bode se graficó cada factor por separado y luego se combinaron gráficamente, debido a los logaritmos implicados. La gráfica de magnitud dependió de la frecuencia, que se midió en rad/s y se ubicó en el eje horizontal logarítmico; el eje vertical presentó la magnitud en decibelios. La gráfica de fase utilizó la misma escala horizontal y presentó la fase en grados.

Un filtro pasivo atenuó o rechazó determinadas frecuencias utilizando resistencias, capacitores e inductores, sin aportar ganancia de potencia. Los filtros activos también utilizaron estos componentes, pero añadieron elementos como amplificadores operacionales, por lo que pudieron ser más fáciles de ajustar y aportar amplificación.

El filtro pasa-altos permitió el paso de frecuencias por encima de una frecuencia de corte \(f_c\) y atenuó las inferiores. El filtro pasa-bajos permitió el paso de frecuencias inferiores a \(f_c\) y atenuó las superiores. El filtro pasa-banda dejó pasar únicamente las frecuencias comprendidas entre una frecuencia inferior \(f_1\) y una superior \(f_2\). El filtro rechaza-banda atenuó las frecuencias comprendidas entre \(f_1\) y \(f_2\) y permitió el paso de las frecuencias externas.

La frecuencia de corte \(f_c\), también denominada frecuencia de \(-3\) dB o de media potencia, fue el valor de frecuencia en el que la amplitud de la señal de salida disminuyó a \(1/\sqrt{2}\approx0.707\) de su valor máximo, equivalente a una caída de 3 dB:\@
\begin{equation}
    A(f_c)=\frac{A_{\max}}{\sqrt{2}}.
    \label{eq:fc}
\end{equation}
Para filtros pasa-banda y rechaza-banda se definieron dos frecuencias de corte, \(f_1\) y \(f_2\), que delimitaron el ancho de banda:
\begin{equation}
    BW=f_2-f_1,\qquad f_0=\sqrt{f_1f_2},
    \label{eq:ancho}
\end{equation}
donde \(BW\) fue el ancho de banda y \(f_0\) la frecuencia central.

\section{Procedimiento experimental}
\subsection{Montaje 1}~\label{mont:1}
\begin{enumerate}
    \item Proceda a montar el circuito de la Figura~\ref{fig:montaje1}. Se utiliza el generador de funciones. Utilice los datos de la Tabla~\ref{tab:montaje1}.
    
    \begin{figure}[htbp]
        \centering
        \begin{circuitikz}[scale=0.8, transform shape]
            \draw(0,0) to[sinusoidal voltage source, l=$V_S$] (0,3);
            \draw(0,3) to[short] (1,3) to[R, l=R] (4,3) to[short] (5,3);
            \draw(5,3) to[C, l=$C$] (5,0);
            \draw(5,0) to[short] (0,0);
            \draw(2.5,0) node[ground]{};
        \end{circuitikz}
        \caption{Montaje 1.}~\label{fig:montaje1}
    \end{figure}

    \begin{table}[H]
        \centering
        \caption{Datos Montaje 1.}~\label{tab:montaje1}
        \begin{tabular}{@{}lc@{}}\toprule
            Voltaje pico & \SI{10}{\volt}\\ 
            Capacitancia & \SI{47}{\nano\farad}\\ 
            Resistencia & \SI{4.7}{\kilo\ohm}\\ 
            Frecuencia 1 & \SI{159}{\hertz}\\ 
            Frecuencia 2 & \SI{1.59}{\kilo\hertz}\\ 
            Frecuencia 3 & \SI{15.9}{\kilo\hertz}\\\bottomrule
        \end{tabular}
    \end{table}

    \item Prepare el Multímetro para realizar mediciones en la señal de entrada y de salida (siendo la señal de salida el voltaje del capacitor).~[\ref{img:1}]
    
    \imagen[img:1]{./Img/img1.png}{Montaje 1.}

    \item Utilizando el multímetro realice las mediciones del voltaje de entrada y de salida con la señal de entrada calibrada al valor de la \textbf{frecuencia 1}.
    
    \begin{flushright}
        Voltaje Eficaz de Entrada:~\SI{7.071}{\volt}\\
        Voltaje Eficaz de Salida:~\SI{6.904}{\volt}
    \end{flushright}

    \item Utilizando el multímetro realice las mediciones del voltaje de entrada y de salida con la señal de entrada calibrada al valor de la \textbf{frecuencia 2}.
    
    \begin{flushright}
        Voltaje Eficaz de Entrada:~\SI{7.071}{\volt}\\
        Voltaje Eficaz de Salida:~\SI{2.911}{\volt}
    \end{flushright}

    \item Utilizando el multímetro realice las mediciones del voltaje de entrada y de salida con la señal de entrada calibrada al valor de la \textbf{frecuencia 3}.
    
    \begin{flushright}
        Voltaje Eficaz de Entrada:~\SI{7.071}{\volt}\\
        Voltaje Eficaz de Salida:~\SI{319.059}{\milli\volt}
    \end{flushright}

    \item Utilizando los valores medidos anteriormente calcule la ganancia de voltaje en decibelios.
    
    \[
        G_{dB} = 20\log_{10}\frac{V_2}{V_1}
    \]

    \begin{flushright}
        Ganancia de Voltaje Frecuencia 1:~\SI{-0.208}{\decibel}\\
        Ganancia de Voltaje Frecuencia 2:~\SI{-7.709}{\decibel}\\
        Ganancia de Voltaje Frecuencia 3:~\SI{-26.912}{\decibel}
    \end{flushright}

    \item Haciendo uso del Bode Plotter obtenga las ganancias de voltaje, según el diagrama de bode del montaje 1 que en este se ilustra, para las frecuencias que se muestran en la Tabla~\ref{tab:montaje1}.
    
    \begin{flushright}
        Ganancia de Voltaje Frecuencia 1:~\SI{-0.209}{\decibel}\\
        Ganancia de Voltaje Frecuencia 2:~\SI{-7.675}{\decibel}\\
        Ganancia de Voltaje Frecuencia 3:~\SI{-26.801}{\decibel}
    \end{flushright}
\end{enumerate}

\subsection{Montaje 2}~\label{mont:2}
\begin{enumerate}
    \item Proceda a montar el circuito de la Figura~\ref{fig:montaje2}. Se utiliza el generador de funciones. Utilice los datos de la Tabla~\ref{tab:montaje2}.
    
    \begin{figure}[htbp]
        \centering
        \begin{circuitikz}[scale=0.8, transform shape]
            \draw(0,0) to[sinusoidal voltage source, l=$V_S$] (0,3);
            \draw(0,3) to[short] (1,3) to[R, l=R] (4,3) to[short] (5,3);
            \draw(5,3) to[L, l=$L$] (5,0);
            \draw(5,0) to[short] (0,0);
            \draw(2.5,0) node[ground]{};
        \end{circuitikz}
        \caption{Montaje 2.}~\label{fig:montaje2}
    \end{figure}
    
    \begin{table}[H]
        \centering
        \caption{Datos Montaje 2.}~\label{tab:montaje2}
        \begin{tabular}{@{}lc@{}}\toprule
            Voltaje pico & \SI{10}{\volt}\\
            Inductancia & \SI{100}{\milli\henry}\\
            Resistencia & \SI{1}{\kilo\ohm}\\
            Frecuencia 1 & \SI{159}{\hertz}\\
            Frecuencia 2 & \SI{1.59}{\kilo\hertz}\\
            Frecuencia 3 & \SI{15.9}{\kilo\hertz}\\\bottomrule
        \end{tabular}
    \end{table}
    
    \item Prepare el Multímetro para realizar mediciones en la señal de entrada y de salida (siendo la señal de salida el voltaje del inductor).
    
    \imagen[fig:2]{./Img/img2.png}{Montaje 2.}
    
    \item Utilizando el multímetro realice las mediciones del voltaje de entrada y de salida con la señal de entrada calibrada al valor de la \textbf{frecuencia 1}.
    
    \begin{flushright}
        Voltaje Eficaz de Entrada:~\SI{7.071}{\volt}\\
        Voltaje Eficaz de Salida:~\SI{705.161}{\milli\volt}
    \end{flushright}

    \item Utilizando el multímetro realice las mediciones del voltaje de entrada y de salida con la señal de entrada calibrada al valor de la \textbf{frecuencia 2}.
    
    \begin{flushright}
        Voltaje Eficaz de Entrada:~\SI{7.071}{\volt}\\
        Voltaje Eficaz de Salida:~\SI{5.006}{\volt}
    \end{flushright}

    \item Utilizando el multímetro realice las mediciones del voltaje de entrada y de salida con la señal de entrada calibrada al valor de la \textbf{frecuencia 3}.
    
    \begin{flushright}
        Voltaje Eficaz de Entrada:~\SI{7.071}{\volt}\\
        Voltaje Eficaz de Salida:~\SI{7.036}{\volt}
    \end{flushright}

    \item Utilizando los valores medidos anteriormente calcule la ganancia de voltaje en decibelios.
    
    \[
        G_{dB} = 20\log_{10}\frac{V_2}{V_1}
    \]

    \begin{flushright}
        Ganancia de Voltaje Frecuencia 1:~\SI{-20.024}{\decibel}\\
        Ganancia de Voltaje Frecuencia 2:~\SI{-3.000}{\decibel}\\
        Ganancia de Voltaje Frecuencia 3:~\SI{-0.043}{\decibel}
    \end{flushright}

    \item Haciendo uso del Bode Plotter obtenga las ganancias de voltaje, según el diagrama de bode del montaje 2 que en este se ilustra, para las frecuencias que se muestran en Tabla~\ref{tab:montaje2}.
    
    \begin{flushright}
        Ganancia de Voltaje Frecuencia 1:~\SI{-19.993}{\decibel}\\
        Ganancia de Voltaje Frecuencia 2:~\SI{-3.073}{\decibel}\\
        Ganancia de Voltaje Frecuencia 3:~\SI{-0.044}{\decibel}
    \end{flushright}
\end{enumerate}

\subsection{Montaje 3}~\label{mont:3}
\begin{enumerate}
    \item Proceda a montar el circuito de la Figura~\ref{fig:montaje3}. Se utiliza el generador de funciones en lugar de una fuente de voltaje (Obligatorio). Utilice los datos de la Tabla~\ref{tab:montaje3}.
    
    \begin{figure}[H]\centering
        \begin{circuitikz}[scale=.8, transform shape]
            \draw(0,0) to[sinusoidal voltage source, l=$V_S$] (0,5) to[short] (1,5) to[R, l=$R_1$] (7,5) to[L, l=$L$] (7,3) to[short] (7,2) to[C, l=$C$] (7,0) to[short] (0,0);
            \draw(5,5) to[R, l=$R_2$] (5,0); \draw(0,0) node[ground]{};
        \end{circuitikz}\caption{Montaje 3.}~\label{fig:montaje3}\end{figure}
    \begin{table}[H]\centering\caption{Datos Montaje 3.}~\label{tab:montaje3}
    
    \begin{tabular}{@{}lc@{}}
        \toprule Voltaje pico & \SI{10}{\volt}\\ 
        Frecuencia & \SI{60}{\hertz}\\ 
        Capacitancia & \SI{33}{\micro\farad}\\ 
        Inductancia & \SI{10}{\milli\henry}\\ 
        Resistencia 1 & \SI{50}{\ohm}\\ 
        Resistencia 2 & \SI{100}{\ohm}\\\bottomrule
    \end{tabular}\end{table}

    \item Ir a Place, luego a Instruments y elegir Bode Plotter. Conectar los terminales: Input (+) al nodo de entrada del circuito Vs, Input (-) a referencia (tierra). Luego Output (+) nodo donde se mide Vo y finalmente Output (-) a referencia (tierra).
    \item Abrir el panel del instrumento y configurar el modo en Magnitude. En el eje horizontal seleccionar barrido Logarítmico (Log) con frecuencia inicial de 1 Hz y frecuencia final de 100 kHz. En el eje vertical seleccionar escala Log (dB) y establecer un rango de -50 dB a 0 dB.
    \item Tome una captura del diagrama identificando con el cursor las frecuencias de corte superior e inferior.
    
    \imagen[fig:3]{./Img/img3.png}{Montaje 3.}
\end{enumerate}

\subsection{Montaje 4}~\label{mont:4}
\begin{enumerate}
    \item Proceda a montar el circuito de la Figura~\ref{fig:montaje4}. Se utiliza el generador de funciones en lugar de una fuente de voltaje (Obligatorio). Utilice los datos de la Tabla~\ref{tab:montaje4}.
    
    \begin{figure}[H]\centering
        \begin{circuitikz}[scale=.8, transform shape]
            \draw(0,0) to[sinusoidal voltage source, l=$V_S$] (0,3) to[short] (1,3) to[R, l=$R_1$] (7,3) to[C, l=$C$] (7,0) to[short] (0,0);
            \draw(5,3) to[L, l=$L$] (5,0); \draw(0,0) node[ground]{};
            \end{circuitikz}
            \caption{Montaje 4. Circuito RLC rechaza-banda.}\label{fig:montaje4}
    \end{figure}
    
    \begin{table}[H]
        \centering
        \caption{Datos Montaje 4.}\label{tab:montaje4}
        \begin{tabular}{@{}lc@{}}
            \toprule 
            Voltaje pico & \SI{10}{\volt}\\ 
            Frecuencia & \SI{60}{\hertz}\\ 
            Capacitancia & \SI{2.2}{\micro\farad}\\ 
            Inductancia & \SI{10}{\milli\henry}\\ 
            Resistencia & \SI{5}{\ohm}\\\bottomrule
        \end{tabular}
    \end{table}
    \item Ir a Place, luego a Instruments y elegir Bode Plotter. Conectar los terminales: Input (+) al nodo de entrada del circuito Vs, Input (-) a referencia (tierra). Luego Output (+) nodo donde se mide Vo y finalmente Output (-) a referencia (tierra).
    \item Abrir el panel del instrumento y configurar el modo en Magnitude. En el eje horizontal seleccionar barrido Logarítmico (Log) con frecuencia inicial de 10 Hz y frecuencia final de 100 kHz. En el eje vertical seleccionar escala Log (dB) y establecer un rango de -30 dB a 5 dB.
    \item Tome una captura del diagrama identificando con el cursor las frecuencias de corte superior e inferior.
    \imagen[fig:4]{./Img/img4.png}{Montaje 4.}
\end{enumerate}

\section{Análisis de resultados}
\begin{enumerate}
    \item Del circuito del montaje~\ref{tab:montaje2}, (Mismos valores de R y L) obtenga la función de transferencia de ganancia de voltaje, esto es: $G=V_o/V_i=V_L/V_S$.
    \begin{quote}
        Al ser un circuito RL en serie, la función de transferencia se obtuvo mediante el divisor de voltaje:
        \[
            V_L = \frac{Ls}{R + Ls}V_S
        \]

        Esto es:
        \[
            G(s) = \frac{V_L}{V_S} = \frac{Ls}{R + Ls}
        \]

        En este caso, tomando en cuenta los valores de $R$ y $L$, obtendríamos:
        \[
            G(s) = \frac{100\times10^{-3}s}{1000 + 100\times10^{-3}s} = \frac{s}{s + 10000}
        \]
    \end{quote}

    \item A partir de la función de transferencia anterior obtenga el diagrama de magnitud y fase haciendo:
    \begin{itemize}
        \item Simplificando la expresión.
        \item Identificando \(K_{dB}\).
        \item Identificando los polos y ceros existentes.
        \item Gráfico de magnitud a partir de los gráficos parciales.
        \item Gráfico de fase a partir de los gráficos parciales.
    \end{itemize}
    \begin{quote}
        La expresión ya la simplificamos en el inciso anterior. Nos queda:
        \[
            G(s) = \frac{s}{s + 10000}
        \]

        La ganancia constante fue \(K_{dB}=0\,\text{dB}\). El sistema tuvo un cero en \(s=0\) y un polo en \(s=-10000\,\si{\radian\per\second}\).

        Como sabemos, $s = j\omega$, donde $\omega = 2\pi f$. Por lo cual se tiene que la amplitud es:

        \[
            |G(f)| = \frac{2\pi f}{\sqrt{10000^2+{(2\pi f)}^2}} = \frac{2\pi f}{\sqrt{100\times10^6 + 4\pi^2 f^2}} = \frac{\pi f}{\sqrt{25\times10^6 + \pi^2 f^2}}
        \]

        Que para el diagrama nos da la magnitud en decibelios como:
        \[
            |G(f)|_{\mathrm{dB}}=20\log_{10}\left(\frac{\pi f}{\sqrt{25\times10^6 + \pi^2 f^2}}\right)
        \]

        En cuanto al ángulo se tiene que:
        \[
            \angle G(f)=\angle(j2\pi f)-\angle(10000+j2\pi f)
            =90^\circ-\tan^{-1}\left(\frac{2\pi f}{10000}\right).
        \]

        \begin{center}
        \begin{tikzpicture}
            \begin{axis}[
                width=0.9\linewidth,
                height=10cm,
                xmode=log,
                xmin=10, xmax=100000,
                ymin=-45, ymax=3,
                grid=both,
                xlabel={Frecuencia (Hz)},
                ylabel={Magnitud (dB)},
                title={Diagrama de magnitud del montaje 2},
                legend style={font=\small, at={(0.02,0.02)}, anchor=south west},
            ]
                \addplot[blue, thick, domain=10:100000, samples=250]
                    {20*ln((2*pi*x*0.1)/sqrt(1000^2+(2*pi*x*0.1)^2))/ln(10)};
                \addlegendentry{Magnitud exacta}
                \addplot[red, dashed, domain=10:1591.549, samples=2]
                    {20*ln((2*pi*x*0.1)/1000)/ln(10)};
                \addlegendentry{Aproximación asintótica}
                \addplot[black, densely dotted] coordinates {(1591.549,-45) (1591.549,3)};
            \end{axis}
        \end{tikzpicture}

        \vspace{0.5em}

        \begin{tikzpicture}
            \begin{axis}[
                width=0.9\linewidth,
                height=6cm,
                xmode=log,
                xmin=10, xmax=100000,
                ymin=0, ymax=100,
                grid=both,
                xlabel={Frecuencia (Hz)},
                ylabel={Fase (grados)},
                title={Diagrama de fase del montaje 2},
                legend style={font=\small, at={(0.02,0.02)}, anchor=south west},
            ]
                \addplot[blue, thick, domain=10:100000, samples=250]
                    {90-atan((2*pi*x)/10000)};
                \addlegendentry{Fase exacta}
                \addplot[black, densely dotted] coordinates {(1591.549,0) (1591.549,100)};
                \addlegendentry{Frecuencia de corte}
            \end{axis}
        \end{tikzpicture}
        \end{center}
    \end{quote}

\item A partir del grafico de ganancia del inciso 2 obtenga la ganancia de voltaje en decibelios a las frecuencias de la tabla 2.
\begin{quote}

    Tomando únicamente los valores de las gráficas, se obtiene lo siguiente:

    \begin{table}[H]
        \centering
        \caption{Ganancia de voltaje obtenida del gráfico de magnitud.}\label{tab:ganancia_grafica_rl}
        \begin{tabular}{@{}lc@{}}
            \toprule
            Frecuencia & Ganancia de voltaje \(G_{dB}\) \\
            \midrule
            \SI{159}{\hertz} & -20\\
            \SI{1.59}{\kilo\hertz} & -3\\
            \SI{15.9}{\kilo\hertz} & 0\\
            \bottomrule
        \end{tabular}
    \end{table}

    De forma aproximada.

\end{quote}

\item Compare los valores de ganancia de voltaje obtenidos en los incisos 6 del procedimiento experimental, con los valores de ganancia de voltajes obtenidos en el inciso 7 del procedimiento experimental. Compare mediante error relativo porcentual. Opine al respecto.
\begin{quote}
    Para realizar la comparación, se hizo uso del error relativo porcentual. Se compara, básicamente, el resultado obtenido por medio de los valores medidos y la fórmulada dada y los valores medidos con el Bode Plotter.
    \begin{equation}
        \text{Error relativo (\%)}=\left|\frac{G_{\mathrm{medida}}-G_{\mathrm{simulada}}}{G_{\mathrm{simulada}}}\right|100.
        \label{eq:error}
    \end{equation}

    Como el inciso no especifíca si es para el montaje 1 o 2, lo haré para ambos colocándo los datos en esta tabla:

\begin{table}[H]
    \centering
    \caption{Comparación entre ganancias medidas y simuladas.}\label{tab:comparacion}
    \begin{tabular}{@{}llccc@{}}
        \toprule
        Montaje & Frecuencia & Medida (dB) & Bode Plotter (dB) & Error (\%) \\
        \midrule
        1 & \SI{159}{\hertz} & \SI{-0.208}{\decibel} & \SI{-0.209}{\decibel} & 0.478\% \\
        1 & \SI{1.59}{\kilo\hertz} & \SI{-7.709}{\decibel} & \SI{-7.675}{\decibel} & 0.443\% \\
        1 & \SI{15.9}{\kilo\hertz} & \SI{-26.912}{\decibel} & \SI{-26.801}{\decibel} & 0.414\% \\
        2 & \SI{159}{\hertz} & \SI{-20.024}{\decibel} & \SI{-19.993}{\decibel} & 0.155\% \\
        2 & \SI{1.59}{\kilo\hertz} & \SI{-3.000}{\decibel} & \SI{-3.073}{\decibel} & 2.376\% \\
        2 & \SI{15.9}{\kilo\hertz} & \SI{-0.043}{\decibel} & \SI{-0.044}{\decibel} & 2.273\% \\
        \bottomrule
    \end{tabular}
\end{table}

\end{quote}

\item A partir de los barridos de frecuencia (gráfico y frecuencia de corte) de los montajes 1, 2, 3 y 4 identifique el tipo de filtro pasivo que representa la configuración.
\begin{quote}
\begin{itemize}
    \item El montaje 1 fue un filtro RC pasa-bajos porque la salida se tomó en el capacitor.
    \item El montaje 2 fue un filtro RL pasa-altos porque la salida se tomó en el inductor.
    \item El montaje 3 fue un filtro RLC pasa-banda, con límites inferior y superior determinados por \(f_1\) y \(f_2\).
    \item El montaje 4 fue un filtro RLC rechaza-banda, con atenuación concentrada entre \(f_1\) y \(f_2\).
\end{itemize}

\end{quote}

\item Investigue la aplicación de cada uno los filtros que identificó en el inciso anterior.
\begin{quote}
El filtro pasa-bajos se utilizó para eliminar ruido de alta frecuencia y suavizar señales. El filtro pasa-altos se aplicó para eliminar componentes de continua y ruido de baja frecuencia, por ejemplo en acoplamiento de audio. El filtro pasa-banda se empleó para seleccionar un canal o una banda específica en sistemas de radio y comunicaciones. El filtro rechaza-banda se aplicó para suprimir una interferencia localizada, como una frecuencia de red, sin eliminar las frecuencias que se encontraban fuera de la banda rechazada.
\end{quote}
\end{enumerate}

\section{Conclusiones}
\begin{itemsinforme}
    \item Los diagramas de Bode resultan ser una forma efectiva y sencilla de analizar el comportamiento de un circuito cuando este es alimentado por frecuencias en un amplio rango. De esta forma, se puede analizar de mejor manera especialmente por su caracter logarítmico que si se analizara con una escala normal.
    \item Durante la práctica se pudieron poner en práctica distintos métodos para el análisis de funciones en transferencia, lo cual nos ayuda a solidificar los conocimientos adquiridos durante la clase.
    \item Vimos, durante esta práctica, lo útil y conveniente que pueden llegar a ser los simuladores para el estudio de circuitos eléctricos. Lo cual fomenta en nosotros el uso de los mismos para distintos contextos.
\end{itemsinforme}

\section{Bibliografía}
\begin{itemize}
    \item Alexander, C. K., Sadiku, M. N. O., y Cordero Pedraza, C. R. (2018). \textit{Fundamentos de circuitos eléctricos} (6.ª ed.). McGraw-Hill/Interamericana.
    \item \textit{Manual de la práctica n.° 4: Diagramas de Bode y filtros}. Laboratorio de Circuitos Eléctricos II, LIE-311.
\end{itemize}

\end{document}
